\chapter{空手接白刃}\label{sec:catching}

\begin{center}
\itshape
庖丁为文惠君解牛，手之所触，肩之所倚，\\
足之所履，膝之所踦，砉然向然，\\
奏刀騞然，莫不中音。\\[4pt]
\upshape ---庄子，《养生主》
\end{center}

\bigskip

Catching inverts manipulation.  In manipulation
(\cref{sec:manipulation}), the robot approaches a stationary
object---the summer insect walks toward the ice.  In catching, the
object approaches the robot: a thrown ball, a tossed tool, a falling
part.  The object is a \emph{sword} in the precise sense of the
framework---it has autonomous actuation (gravity plus initial
velocity), it is observable (camera), and its trajectory intersects
the king's reachable set.  Catching is path~(a) resolution: upon
capture, the sword relinquishes autonomy ($U_r \to \varnothing$,
velocity $\to 0$, object at rest in the gripper).

The key insight is that the catcher's technique is invariant.
The equivalence principle of general relativity---inertial mass
equals gravitational mass---implies that the object's free-flight
trajectory does not depend on its mass.  The controller operates in
the tangent bundle $TM$, where mass has already cancelled.  The
catcher does not need to know \emph{what} it is catching.  Like
庖丁解牛, the blade follows the structure; the hand follows the
gap.

The argument proceeds in seven stages:
\begin{enumerate}
  \item \textbf{Diagnosis} (\cref{sec:c-diagnosis}): the approaching
  sword---autonomous actuation via Newtonian mechanics.
  \item \textbf{Equivalence principle} (\cref{sec:c-equivalence}):
  mass cancellation in $TM$---the controller is universal.
  \item \textbf{Potential well} (\cref{sec:c-well}): chopstick
  mechanics---two elastic rods create a harmonic trap.
  \item \textbf{Capture theorem} (\cref{sec:c-capture}): the
  topological transition $\lambda_1\colon 0 \to 3$.
  \item \textbf{Three-term decomposition} (\cref{sec:c-three-term}):
  predict $+$ intercept $+$ bang.
  \item \textbf{Kakeya duality} (\cref{sec:c-kakeya}): catching
  $\leftrightarrow$ maintenance under time reversal.
  \item \textbf{Implementation} (\cref{sec:c-impl}): MuJoCo
  simulation and Flexiv Rizon~4 hardware.
\end{enumerate}


% ═══════════════════════════════════════════════════════════
\section{Diagnosis: the approaching sword}
\label{sec:c-diagnosis}

\subsection{The sword in free flight}

\begin{definition}[Approaching sword]
\label{def:c-sword}
An object $r$ in free flight is an \emph{approaching sword} if
three conditions hold:
\begin{enumerate}[label=(\roman*),nosep]
  \item $U_r \neq \varnothing$: the object has autonomous actuation
  via gravity and initial velocity $v_0$;
  \item $r \in \mathrm{Im}(\Obs)$: the object is observable by the
  king's sensor (camera);
  \item $\exists\, t^* < \infty$: the object's trajectory intersects
  the king's reachable set $\mathrm{Reach}(\kappa)$.
\end{enumerate}
\end{definition}

The approaching sword is the purest instantiation of the framework's
sword concept.  In the main text, ``autonomous actuation'' is
abstract---an adversary, a competitor, a market force.  Here it is
literal physics: $F = mg$.  The actuation is Newtonian mechanics
itself~\cite{newton}.

\begin{proposition}[Sword lifecycle of a thrown object]
\label{prop:c-lifecycle}
The approaching sword resolves via path~(a) of the binary fate
theorem (\cref{thm:lifecycle}): upon capture, $U_r \to \varnothing$
(velocity $\to 0$, object at rest in the gripper).  Path~(b)
(elimination) corresponds to deflection or dodging.  No path~(c)
exists---the object cannot remain in flight indefinitely within
$\mathrm{Reach}(\kappa)$.
\end{proposition}

\begin{proof}
The object follows the parabolic trajectory
$x(t) = x_0 + v_0 t + \frac{1}{2}g t^2$.  Since $\|g\| > 0$,
the trajectory is unbounded: $\|x(t)\| \to \infty$ as
$t \to \infty$.  Therefore the object either intersects
$\mathrm{Reach}(\kappa)$ at finite $t^*$ (capture or deflection)
or exits the workspace permanently.  In either case, the sword
does not persist.
\end{proof}

\begin{remark}[Connection to the forcing lemma]
The forcing lemma (\cref{lem:forcing}) states that a persistent
sword forces the king to act.  Here persistence is bounded by
physics: gravity ensures that no thrown object persists.  The
approaching sword is a sword with a built-in deadline $t^*$.
The catching problem is therefore time-constrained in a way that
manipulation is not---the summer insect can take its time, but
the catcher cannot.
\end{remark}

\subsection{Mapping table}

\begin{center}
\begin{tabular}{@{}lll@{}}
\toprule
\textbf{Catching} & \textbf{Framework} & \textbf{Manipulation} \\
\midrule
Thrown object & Sword & Stationary object (passive) \\
Gravity $+ v_0$ & Autonomous actuation $U_r$ & None (object is passive) \\
Camera & Detection $\Obs$ & Wrist camera \\
Catching & Path~(a): relinquish & Phase~2: loading \\
Dodging & Path~(b): elimination & N/A \\
Deadline $t^*$ & Forcing & No deadline \\
\bottomrule
\end{tabular}
\end{center}


% ═══════════════════════════════════════════════════════════
\section{The equivalence principle}
\label{sec:c-equivalence}

\subsection{Mass-free trajectory}

\begin{proposition}[Mass-free trajectory]
\label{prop:c-massfree}
The trajectory of a thrown object in free flight,
\[
  \ddot{x}(t) = g,\qquad
  x(0) = x_0,\quad \dot{x}(0) = v_0,
\]
is independent of mass~$m$.  The intercept point
$x^* = x(t^*)$ and intercept velocity $v^* = \dot{x}(t^*)$
depend only on the initial conditions $(x_0, v_0)$ and the
gravitational acceleration~$g$.
\end{proposition}

\begin{proof}
Newton's second law gives $ma = mg$, hence $a = g$
for all $m > 0$.  The trajectory
$x(t) = x_0 + v_0 t + \frac{1}{2}g t^2$ and velocity
$\dot{x}(t) = v_0 + g t$ contain no $m$.
\end{proof}

\begin{remark}[Galileo's tower]
This is Galileo's insight at Pisa: all objects fall at the
same rate.  In general relativity, it becomes the equivalence
principle---inertial mass equals gravitational mass.  In our
context: the robot's perception pipeline (camera $\to$
position $\to$ velocity $\to$ parabolic fit) operates
entirely in configuration space, where mass has already
cancelled.  The robot does not need to know what it is
catching.
\end{remark}

\subsection{The mass-invariant controller}

\begin{theorem}[Mass-invariant catching]
\label{thm:c-equivalence}
Let $\pi\colon (x_{\mathrm{obj}}, \dot{x}_{\mathrm{obj}},
q, \dot{q}) \mapsto \tau$ be a catching controller that operates
in the tangent bundle $TM$ of the configuration manifold.
Then the control policy $\pi$ is invariant under mass rescaling
$m \to \alpha m$ for all $\alpha > 0$, provided the end-effector
satisfies the hardware envelope $K > mv_\perp^2/d^2$.
\end{theorem}

\begin{proof}
We verify mass cancellation in each of the three stages of
catching.

\emph{Stage~1: Pre-contact (trajectory prediction).}
By \cref{prop:c-massfree}, the free-flight trajectory
$x(t) = x_0 + v_0 t + \frac{1}{2}g t^2$ is mass-free.
The intercept point $x^*$ and intercept time $t^*$ are
computed from $(x_0, v_0, g)$ alone.

\emph{Stage~2: Contact (impedance response).}
At contact, the chopstick pair applies impedance control:
\[
  \tau = K(x_d - x) + D(\dot{x}_d - \dot{x}),
\]
where $K$ is the spring stiffness and $D$ is the damping
coefficient.  Both are properties of the end-effector, not
of the object.  The physical spring absorbs kinetic energy
$E = \frac{1}{2}mv^2$ via compression
$\delta = v\sqrt{m/K}$, but the controller does not need
to know $\delta$---the spring self-adjusts.

\emph{Stage~3: Post-contact (static hold).}
The grasp quality metric
$\rho = \|F_{\mathrm{contact}}\| / (mg)$ has both
numerator and denominator proportional to $m$, so $\rho$
is mass-invariant.
\end{proof}

\begin{remark}[Hardware envelope]
\label{rem:c-envelope}
The policy $\pi$ is universal within the \emph{hardware design
envelope}: the set of masses $m$ satisfying $K > mv_\perp^2/d^2$
(\cref{prop:c-trap}).  The viability window
$\Delta t_{\mathrm{viable}}$ also depends on $m$ through the
decay rate $m/(2D)$.  Universality is in the software: the same
policy catches tennis balls and bowling balls; only the spring
stiffness differs.
\end{remark}

\begin{remark}[Connection to the mass gap]
In the mass-gap theorem (\cref{thm:massgap}), the positivity
$\lambda_1 > 0$ is equivalent to graph connectivity---it does
not depend on whether edge weights are large or small, only on
whether they are positive.  Here, the catching policy does not
depend on whether the object is heavy or light, only on whether
it enters the well.  Both are instances of universality from
connectivity: the qualitative structure (connected vs.\
disconnected, caught vs.\ uncaught) is robust; the quantitative
details (specific mass, specific edge weight) determine margins
but not the outcome.
\end{remark}


% ═══════════════════════════════════════════════════════════
\section{The potential well}
\label{sec:c-well}

\subsection{Chopstick pair}

\begin{definition}[Chopstick pair]
\label{def:c-chopstick}
A \emph{chopstick pair} $(C_L, C_R)$ consists of two elastic
rods, each with:
\begin{enumerate}[label=(\roman*),nosep]
  \item rigid base attachment to a gripper finger (no additional
  degree of freedom);
  \item a spring-loaded hinge at position $s \in (0, L)$ with
  stiffness $K$ and damping coefficient $D$;
  \item a contact tip at position $s = L$.
\end{enumerate}
When the chopstick pair is closed to finger separation~$2d$,
the pair creates a potential field
\[
  V_{\mathrm{trap}}(x)
  = \tfrac{1}{2}K_L\,(x - x_L)^2
  + \tfrac{1}{2}K_R\,(x - x_R)^2
\]
with minimum at the midpoint
$\bar{x} = \frac{x_L + x_R}{2}$.
\end{definition}

The chopstick pair is the physical instantiation of the
potential well.  Two compliant fingers, each with a spring
hinge, create a region of space where incoming kinetic energy
is converted to potential energy---the object decelerates,
reverses, and oscillates until damping brings it to rest.

\subsection{Harmonic trapping}

\begin{proposition}[Harmonic trapping]
\label{prop:c-trap}
An object entering the chopstick pair with transverse velocity~$v$
undergoes damped oscillation:
\[
  m\ddot{x} + 2D\dot{x} + 2Kx = 0,
\]
with characteristic frequency $\omega = \sqrt{2K/m}$ and decay
time $\tau_d = m/(2D)$.

The object is trapped (velocity reversal within the well) if
and only if
\[
  \tfrac{1}{2}mv^2 < V_{\mathrm{trap}}(d)
  = \tfrac{1}{2}K d^2,
\]
equivalently, $v < d\sqrt{K/m}$.
\end{proposition}

\begin{proof}
The equation of motion follows from superposing the two spring
forces: $F = -K_L(x - x_L) - K_R(x - x_R) - 2D\dot{x}$.
With $K_L = K_R = K$ and origin at $\bar{x}$, this reduces to
$m\ddot{x} + 2D\dot{x} + 2Kx = 0$.  The trapping condition
is energy conservation: the object's kinetic energy at entry
must not exceed the potential barrier at the well boundary~$d$.
\end{proof}

\begin{remark}[Mass in the spring, not in the controller]
The trapping condition $v < d\sqrt{K/m}$ contains~$m$, but
this is a constraint on the \emph{spring design}, not on the
\emph{controller}.  The controller's task is to deliver the
chopstick tips to the intercept point---the spring does the
rest.  For a given stiffness~$K$, heavier objects with the
same velocity require stiffer springs (larger~$K$), but the
control policy~$\pi$ is identical.  Mass enters the hardware
specification, not the software.
\end{remark}

\subsection{Potential well diagram}

\begin{center}
\begin{tikzpicture}[>=stealth, scale=0.9]
  % Potential well curve
  \draw[thick, water] (-3,4) parabola bend (0,0.5) (3,4);
  \node[water, above] at (0,0.3) {$\bar{x}$};

  % Well boundary markers
  \draw[dashed, gray] (-2,0) -- (-2,4.2);
  \draw[dashed, gray] (2,0) -- (2,4.2);
  \node[below] at (-2,0) {$x_L$};
  \node[below] at (2,0) {$x_R$};

  % Separation marker
  \draw[<->] (-2,-0.4) -- (2,-0.4);
  \node[below] at (0,-0.4) {$2d$};

  % Springs (left and right)
  \draw[thick, dao, decorate,
    decoration={coil, segment length=4pt, amplitude=3pt}]
    (-3.2,2) -- (-2,2);
  \draw[thick, dao, decorate,
    decoration={coil, segment length=4pt, amplitude=3pt}]
    (2,2) -- (3.2,2);
  \node[dao, left] at (-3.2,2) {$K,D$};
  \node[dao, right] at (3.2,2) {$K,D$};

  % Gripper fingers
  \fill[gray!30] (-3.8,0) rectangle (-3.2,4);
  \fill[gray!30] (3.2,0) rectangle (3.8,4);
  \node[rotate=90] at (-3.5,2) {\small finger};
  \node[rotate=90] at (3.5,2) {\small finger};

  % Object trajectory (entering from above)
  \draw[thick, sword, ->] (0,5.5) -- (0,4.2);
  \filldraw[sword] (0,4.2) circle (4pt);
  \node[sword, right] at (0.3,5) {$v_\perp$};

  % Energy level
  \draw[thin, caution, dashed] (-2.5,3.2) -- (2.5,3.2);
  \node[caution, right] at (2.5,3.2)
    {$\frac{1}{2}mv_\perp^2$};

  % Potential label
  \node[water, left] at (-3,4) {$V_{\mathrm{trap}}(x)$};

  % Axes
  \draw[->] (-4,0) -- (4,0) node[right] {$x$};
  \draw[->] (0,0) -- (0,6) node[above] {$V$};
\end{tikzpicture}
\end{center}


% ═══════════════════════════════════════════════════════════
\section{The capture theorem}
\label{sec:c-capture}

This section contains the core theorem of the chapter: the
\emph{potential well capture theorem}, which characterises the
topological transition from $\lambda_1 = 0$ (no contact edges)
to $\lambda_1 > 0$ (closed contact cycle).

\subsection{Viability window}

\begin{definition}[Viability window]
\label{def:c-viability-window}
The \emph{viability window} $\Delta t_{\mathrm{viable}}$ is the
time interval during which the gripper can close to seal the
potential well:
\[
  \Delta t_{\mathrm{viable}}
  = \frac{x_{\mathrm{escape}}}{\dot{x}_{\mathrm{residual}}},
\]
where $x_{\mathrm{escape}}$ is the distance from the object to the
open end of the well and $\dot{x}_{\mathrm{residual}}$ is the
object's residual velocity after partial spring absorption.
\end{definition}

Here $\Delta t_{\mathrm{absorb}}$ denotes the elapsed time
between the object entering the chopstick well (first spring
contact) and the gripper beginning to close.  During this
interval the springs decelerate the object; the gripper
remains open.

The viability window is the catching analogue of the forcing
lemma's deadline: the catcher must act before the object escapes
the open well.  After the springs absorb part of the kinetic
energy, the object still moves---the well is open at one end.
The gripper must close before the object exits.

\subsection{The theorem}

\begin{theorem}[Potential well capture --- 势井捕获定理]
\label{thm:c-capture}
Let $(C_L, C_R)$ be a chopstick pair with spring stiffness~$K$
and damping coefficient~$D$, mounted on a gripper with closure
time $\Delta t_{\mathrm{grip}}$.  An approaching sword with
transverse velocity $v_\perp$ at the intercept point is captured
($\lambda_1\colon 0 \to 3$) if and only if three conditions
hold:
\begin{enumerate}[label=(\roman*),nosep]
  \item \textbf{Well formation:} the arm delivers the chopstick
  tips to the intercept point before the object arrives:
  $t_{\mathrm{arm}} < t^*$.
  \item \textbf{Energy absorption:} the transverse kinetic energy
  is within the spring capacity:
  $\frac{1}{2}mv_\perp^2 < \frac{1}{2}Kd^2$.
  \item \textbf{Timing:} the gripper closes within the viability
  window:
  $\Delta t_{\mathrm{grip}} < \Delta t_{\mathrm{viable}}$.
\end{enumerate}
Under these conditions, the contact graph undergoes the transition
\[
  \bar{K}_3 \;\text{(three isolated nodes)}
  \xrightarrow{\;\text{gripper bang}\;}
  C_3 \;\text{(cycle graph)},
\]
where $\lambda_1(\bar{K}_3) = 0$ (the graph is disconnected:
the gripper is open and no rigid contact edges exist) and
$\lambda_1(C_3) = 3 > 0$ (the cycle has maximal spectral gap
for three nodes).  During the viability window, the springs
decelerate the object but do not constitute rigid contact
edges in the sense of the graph Laplacian
(\cref{thm:massgap}).
\end{theorem}

\begin{proof}
We verify each condition and the resulting topological transition.

\emph{Condition~(i): well formation.}
This is a reachability condition on the robot arm.  Given the
intercept point $x^* = x_0 + v_0 t^* + \frac{1}{2}g (t^*)^2$
computed from \cref{prop:c-massfree}, the arm must solve the
inverse kinematics $q^* = \mathrm{IK}(x^*)$ and reach $q^*$
within time~$t^*$.  This is feasible whenever
$x^* \in \mathrm{Reach}(\kappa)$ and the arm's maximum
joint velocity exceeds the required trajectory.

\emph{Condition~(ii): energy absorption.}
By \cref{prop:c-trap}, the chopstick pair traps the object if
$\frac{1}{2}mv_\perp^2 < \frac{1}{2}Kd^2$.  After entering
the well, the object undergoes damped oscillation with
residual velocity
$\dot{x}_{\mathrm{res}}(t) = v_\perp e^{-Dt/m}$.
The spring absorbs kinetic energy continuously; total absorption
occurs as $t \to \infty$, but significant deceleration occurs
within $\tau_d = m/(2D)$.

\emph{Condition~(iii): timing.}
After the object enters the well, its residual velocity is
$\dot{x}_{\mathrm{res}} = v_\perp e^{-D\Delta t_{\mathrm{absorb}}/m}$
at the moment the gripper begins closing.  The escape distance
is $x_{\mathrm{escape}} \sim L_{\mathrm{chopstick}}$.
The viability window is therefore
\[
  \Delta t_{\mathrm{viable}}
  = \frac{L}{\dot{x}_{\mathrm{res}}}
  = \frac{L}{v_\perp}\, e^{D\Delta t_{\mathrm{absorb}}/m}.
\]
For representative parameters ($L \sim 0.1$\,m,
$v_\perp \sim 1$\,m/s post-absorption),
$\Delta t_{\mathrm{viable}} \sim 100$\,ms, which exceeds
typical gripper closure times
$\Delta t_{\mathrm{grip}} \sim 50$\,ms.

\emph{Topological transition.}
Before gripper closure, the contact graph has three nodes
$\{C_L, \mathrm{obj}, C_R\}$ and \emph{no edges}: the gripper
is open, so neither finger forms a rigid contact with the object,
and the fingers are not rigidly connected to each other.  The
graph is the independent set~$\bar{K}_3$, whose Laplacian is the
zero matrix; in particular $\lambda_1 = 0$
(\cref{thm:massgap}: disconnected $\Leftrightarrow$
$\lambda_1 = 0$).  During the viability window, the chopstick
springs decelerate the object through compliant forces, but
compliant contact does not constitute a rigid edge.

When the gripper fires (the bang arc), all three rigid contact
edges form simultaneously:
$C_L\text{---}\mathrm{obj}$,
$\mathrm{obj}\text{---}C_R$, and
$C_L\text{---}C_R$ (through the gripper body).
The resulting cycle~$C_3$ has Laplacian eigenvalues
$\{0, 3, 3\}$, so $\lambda_1(C_3) = 3 > 0$~\cite{cheeger}.
The spectral gap jumps from $0$ to~$3$ in a single discrete
event: the object is enclosed.
\end{proof}

\begin{remark}[The bang is the spectral kick]
The gripper closure is Term~(III) of the three-term
decomposition---the bang arc.  It is the discrete topological
event that creates $\lambda_1 = 3$ from $\lambda_1 = 0$---three isolated nodes become a cycle.
All of catching reduces to delivering Term~(III) at the right
moment within the viability window.  The singular arc
(Term~II, arm trajectory) is smooth and continuous; the bang
arc is discrete and instantaneous.  The spectral gap is born
in that instant.
\end{remark}

\begin{remark}[空手接白刃]
The name 空手接白刃 means ``catching a blade bare-handed.''
In martial arts, this requires catching the blade between two
palms at the exact moment of minimal relative velocity---after
the swing's kinetic energy is partially absorbed by the body's
compliance, and before the attacker can retract.  This is
precisely the viability window: the interval between spring
absorption and object escape.  The martial artist's two palms
are the chopstick pair; the clap is the gripper bang.
庖丁's blade finds the gap between sinew and bone; the
catcher's hands find the gap between arrival and escape.
\end{remark}


% ═══════════════════════════════════════════════════════════
\section{Three-term decomposition}
\label{sec:c-three-term}

\begin{theorem}[Catching decomposition]
\label{thm:c-decomp}
The optimal catching control decomposes into three terms:
\[
  \tau^*(t)
  = \underbrace{\tau_{\mathrm{predict}}(t)}_{\text{(I)\;gravity/drift}}
  + \underbrace{\tau_{\mathrm{intercept}}(t)}_{\text{(II)\;least action}}
  + \underbrace{\tau_{\mathrm{bang}}(t)}_{\text{(III)\;spectral kick}},
\]
where:
\begin{enumerate}[label=(\Roman*),nosep]
  \item \textbf{Prediction:} analytical solution of the projectile
  equation $\ddot{x} = g$ yields the intercept point $(x^*, t^*)$.
  Zero control effort---computation only.
  \item \textbf{Intercept (singular arc):} minimum-jerk trajectory
  of the arm to~$x^*$.  Occupies approximately $90\%$ of catching
  duration.  Smooth, differentiable.
  \item \textbf{Gripper bang (bang arc):} close gripper at
  $t = t^* + \Delta t_{\mathrm{absorb}}$.  Discrete, topological.
  Occupies approximately $10\%$ of duration.
\end{enumerate}
\end{theorem}

\begin{proof}
By the Pontryagin maximum principle, the optimal control
switches between singular and bang arcs at the switching surface
$\Sigma = \{\rho = 1\}$.  The Hamiltonian
\[
  H = p^\top f(x,u) + \lambda_0 L(x,u)
\]
is linear in $u$ on the bang arc (gripper: open/closed) and
satisfies $\partial H / \partial u = 0$ on the singular arc
(arm trajectory).  This is the same structure as the gravity
damper (\cref{thm:3b-damper}): Term~(I) is the
uncontrolled drift, Term~(II) is the smooth tracking, and
Term~(III) is the discrete impulse at $\Sigma$.
\end{proof}

\subsection{Mapping table}

\begin{center}
\begin{tabular}{@{}lllll@{}}
\toprule
\textbf{Term} & \textbf{Three-body} & \textbf{Manipulation}
  & \textbf{Catching} \\
\midrule
(I) Gravity/drift
  & Gravitational field
  & Object weight
  & Projectile trajectory \\
(II) Least action
  & Damper tracks apex
  & Arm approaches object
  & Arm moves to intercept \\
(III) Spectral kick
  & Damper impulse
  & Gripper close (load)
  & Gripper bang (capture) \\
Switching $\Sigma$
  & $\rho = 1$ contact
  & $\rho = 1$ grasp
  & $\rho = 1$ catch \\
Time fraction (III)
  & ${\sim}10\%$
  & ${\sim}10\%$
  & ${\sim}10\%$ \\
\bottomrule
\end{tabular}
\end{center}

\begin{remark}[Universality of the $10\%$ ratio]
Across all three instantiations---celestial mechanics,
tabletop manipulation, and projectile catching---the bang arc
occupies approximately $10\%$ of the total control duration.
This is an empirical regularity: the bang arc is the
topological event (contact creation or destruction), which
occurs at the codimension-$1$ switching surface $\Sigma$,
while the singular arc is the smooth approach through the
interior.  The ${\sim}10\%$ ratio reflects the geometric
brevity of the transition relative to the approach, though
the exact fraction depends on hardware timescales (gripper
closure speed, actuator bandwidth).
\end{remark}


% ═══════════════════════════════════════════════════════════
\section{Kakeya duality}
\label{sec:c-kakeya}

\begin{theorem}[Catching--maintenance duality]
\label{thm:c-kakeya}
Let $\mathcal{M}$ denote the maintenance problem (keep
$\lambda_1 > 0$, prevent escape from the potential well) and
$\mathcal{C}$ denote the capture problem (achieve $\lambda_1 > 0$,
trap the object in the potential well).  Then $\mathcal{C}$ and
$\mathcal{M}$ are related by time reversal:
\[
  \mathcal{C} = \mathcal{T} \circ \mathcal{M}
  \qquad \text{(exact when $D = 0$)},
\]
where $\mathcal{T}\colon t \mapsto -t$ is the time-reversal
operator.  With dissipation ($D > 0$), the duality is
approximate at the trajectory level but exact at the
topological level (the contact graph transition is identical).
\end{theorem}

\begin{proof}
Consider the time-reversed maintenance problem.  In maintenance,
the object is inside the well at $t = 0$ with $\lambda_1 > 0$,
and the controller prevents escape for all $t > 0$.  Under time
reversal $t \mapsto -t$, the object starts outside the well at
$t = 0$ with $\lambda_1 = 0$, and the controller must bring it
inside by $t = t^*$ with $\lambda_1 > 0$.  This is precisely the
capture problem.

The duality is exact in the conservative case ($D = 0$):
reversing the velocity trajectory of a successful maintenance
episode yields a valid capture trajectory.  With dissipation
($D > 0$), time reversal is approximate---the reversed trajectory
gains energy rather than losing it---but the topological structure
(which nodes are connected) is identical.
\end{proof}

\begin{center}
\begin{tabular}{@{}lll@{}}
\toprule
\textbf{Property} & \textbf{Maintenance $\mathcal{M}$}
  & \textbf{Capture $\mathcal{C}$} \\
\midrule
Initial $\lambda_1$ & $> 0$ & $= 0$ \\
Target $\lambda_1$ & $> 0$ (maintain) & $> 0$ (achieve) \\
Potential well & Exists; prevent escape
  & Does not exist; must create \\
Direction & Inside $\to$ stay inside
  & Outside $\to$ bring inside \\
Kakeya condition & Free stability impossible
  & Free capture impossible \\
\bottomrule
\end{tabular}
\end{center}

\begin{corollary}[Kakeya condition for catching]
\label{cor:c-kakeya}
There exists no passive mechanism that captures all approaching
swords without active control.  The arm must actively position
the chopstick pair (Term~II) and actively close the gripper
(Term~III).  The springs provide passive \emph{absorption} but
not passive \emph{capture}.
\end{corollary}

\begin{proof}
Suppose a passive mechanism (no active control, $\tau = 0$)
could capture all approaching swords.  Then the potential well
would need to be pre-positioned at all possible intercept
points simultaneously.  But the set of intercept points
$\{x^* : (x_0, v_0) \in \mathcal{S}\}$ for the full set of
initial conditions~$\mathcal{S}$ has positive measure in
$\R^3$.  A single passive well occupies zero measure.
Therefore no passive mechanism suffices.

This is the catching analogue of the Kakeya condition: free
stability does not exist, free capture does not exist.
Active control is necessary.
\end{proof}

\begin{remark}[The trap must come to the sword]
The Kakeya duality reveals the deep structure: a trap
placed at a fixed location catches only objects that happen
to arrive there.  A universal catcher must \emph{move} the
trap to the predicted intercept point.  This is why
Term~(II) exists: the singular arc is the transport of the
potential well through space.  The well is not a place---it
is an event: the right configuration at the right time.
\end{remark}


% ═══════════════════════════════════════════════════════════
\section{Implementation: Flexiv Rizon~4}
\label{sec:c-impl}

\subsection{Hardware platform}

The Flexiv Rizon~4 provides the physical substrate:
\begin{itemize}[nosep]
  \item 7 degrees of freedom (anthropomorphic arm);
  \item 1\,kHz real-time control loop (5\,kHz internal servo);
  \item 0.03\,N force sensing resolution at end-effector;
  \item 4\,kg payload capacity;
  \item joint-level torque control interface.
\end{itemize}

\subsection{Gripper dynamics: bang-bang contact}

The Flexiv gripper is a parallel jaw mechanism driven by a
gear-and-rack transmission.  Its control is bang-bang:
\begin{enumerate}[nosep]
  \item \textbf{Free space} (gripper open): position control
  drives the jaws to the pre-set aperture.  No contact forces.
  \item \textbf{Contact ramp-up} (jaws reach object): motor
  current rises until the force threshold $F_{\mathrm{thresh}}$
  is reached.  The contact dynamics satisfy the Signorini
  complementarity condition:
  \[
    d \geq 0,\quad F \geq 0,\quad d \cdot F = 0,
  \]
  where $d$ is the jaw--object gap and $F$ is the contact force.
  At contact ($d = 0$), the force ramps to threshold.
  \item \textbf{Static hold} (mechanical self-lock): the
  transmission's worm gear is non-backdrivable---once the
  threshold is reached, the jaws lock mechanically without
  active motor torque.  The contact graph $C_3$ persists
  with $\lambda_1 = 3$ at zero energy cost.
\end{enumerate}
The bang-bang structure is not a simplification---it is the
correct model.  The switching surface $\Sigma = \{\rho = 1\}$
corresponds to $d = 0$ (contact established).  The force
threshold $F_{\mathrm{thresh}}$ is a Signorini constraint,
not a PID setpoint.  The mechanical self-lock is the physical
realisation of spectral gap maintenance: once $\lambda_1 > 0$,
it persists without control effort.

\begin{remark}[Why not impedance control?]
Impedance control ($F = M\ddot{x} + B\dot{x} + K(x - x_d)$)
requires high-bandwidth force sensing and backdrivable
actuators.  Industrial parallel grippers have neither: the
worm gear introduces high friction and non-backdrivability.
For catching, this is an advantage: the non-backdrivability
\emph{is} the mechanical self-lock that maintains $\lambda_1 > 0$
after capture.  The compliance needed for energy absorption
lives in the chopstick springs, not in the gripper actuator.
The separation is clean: springs absorb (passive compliance),
gripper locks (active topology change).
\end{remark}

\subsection{Chopstick end-effector: X-geometry}

The chopstick pair is realised as two carbon-fibre rods
($L = 150$\,mm, $\varnothing = 3$\,mm) mounted in an
X-crossing geometry:
\begin{itemize}[nosep]
  \item the rods cross at their midpoints, forming an X;
  \item a torsion spring at the crossing point
  ($K = 0.5$\,N$\cdot$m/rad,
  $D = 0.02$\,N$\cdot$m$\cdot$s/rad) provides compliance;
  \item proximal tips (above crossing) attach to the gripper
  fingers;
  \item distal tips (below crossing) form the catching well,
  with silicone contact pads ($\mu \geq 0.8$).
\end{itemize}
The X-geometry converts parallel finger closure (at the
proximal tips) into well closure at the distal tips via the
crossing lever arm.  When the gripper is open, the distal
tips spread apart, creating a funnel that guides the
incoming object toward the spring.  When the gripper closes
(bang), the distal tips converge, sealing the well.
The resulting potential well has depth
$V_{\mathrm{trap}}(d) = \frac{1}{2}Kd^2 \approx 0.3$\,J
for $d = 35$\,mm, sufficient to capture objects up to
$v_\perp \approx 2$\,m/s at $m = 0.15$\,kg.

\subsection{Catching action bound}

\begin{proposition}[Catching action bound]
\label{prop:c-action-bound}
The hardware configuration
$(K, D, d, L, \Delta t_{\mathrm{grip}})$ determines the
catching envelope.  Define the \emph{energy margin}
$\eta_E = Kd^2/(mv_\perp^2)$ and the \emph{timing margin}
$\eta_T = \Delta t_{\mathrm{viable}} /
\Delta t_{\mathrm{grip}}$.  An approaching sword with
mass~$m$ and transverse velocity~$v_\perp$ is capturable if
and only if both margins exceed unity:
\[
  \eta_E > 1 \quad\text{and}\quad \eta_T > 1.
\]
The product $\eta = \eta_E \cdot \eta_T$ is the
\emph{catching number}; $\eta > 1$ is necessary but the
binding constraint is $\min(\eta_E, \eta_T) > 1$.  The
boundary where the tighter margin equals unity is the
\emph{action extremum}: the catch that exactly saturates
the hardware budget.
\end{proposition}

\begin{proof}
The trapping condition (\cref{prop:c-trap}) requires
$\frac{1}{2}mv_\perp^2 < \frac{1}{2}Kd^2$, i.e.\
$Kd^2/(mv_\perp^2) > 1$.  The timing condition
(\cref{thm:c-capture}(iii)) requires
$\Delta t_{\mathrm{grip}} < \Delta t_{\mathrm{viable}}$.
Capture requires both margins individually: $\eta_E > 1$
(enough spring to trap) and $\eta_T > 1$ (enough time to
close).  The product $\eta = \eta_E \cdot \eta_T > 1$ is
necessary but not sufficient; the binding constraint is
the smaller margin.

For the Flexiv configuration ($K_{\mathrm{eff}} \approx 490$\,N/m,
$d = 35$\,mm, $\Delta t_{\mathrm{grip}} = 50$\,ms,
$L = 150$\,mm), the boundary $\eta = 1$ traces a curve in the
$(m, v_\perp)$ plane: heavier objects require slower approach
(stiffer springs widen the envelope), faster objects require
lighter mass (longer viability window).
\end{proof}

\begin{remark}[Action as energy $\times$ time]
The two margins decompose as
\[
  \eta_E = \frac{E_{\mathrm{trap}}}{E_{\mathrm{kin}}},
  \qquad
  \eta_T = \frac{\Delta t_{\mathrm{viable}}}
  {\Delta t_{\mathrm{grip}}},
\]
where $E_{\mathrm{trap}} = \frac{1}{2}Kd^2$ is the well depth
and $E_{\mathrm{kin}} = \frac{1}{2}mv_\perp^2$ is the object's
kinetic energy.  $\eta_E$ is the energy action ratio (can the
spring absorb the kinetic energy?), $\eta_T$ is the timing
action ratio (can the gripper close before the object escapes?).
The catch succeeds when \emph{both} exceed unity.  The
optimal catch is the one that saturates the tighter margin:
$\min(\eta_E, \eta_T) = 1$.
\end{remark}

\subsection{MuJoCo simulation}

The simulation environment models the chopstick mechanics using
MuJoCo's~\cite{mujoco} native spring-damper system:
\begin{itemize}[nosep]
  \item hinge joints with \texttt{stiffness} and \texttt{damping}
  attributes matching the physical springs;
  \item \texttt{springref} set to the closed-chopstick angle;
  \item condim~4 contact model for frictional grasping;
  \item 2\,ms simulation timestep (500\,Hz).
\end{itemize}

\subsection{Perception pipeline}

The catching pipeline operates in four stages:
\begin{enumerate}[nosep]
  \item RGB-D image $\to$ colour filter $\to$ object
  segmentation;
  \item depth map $\to$ 3D position $x(t)$;
  \item two-frame finite difference $\to$ velocity
  $\dot{x}(t) \approx (x(t) - x(t - \Delta t))/\Delta t$;
  \item parabolic fit $\to$ intercept prediction $(x^*, t^*)$
  via \cref{prop:c-massfree}.
\end{enumerate}
The full pipeline runs at 30\,Hz (camera frame rate).
Prediction latency from image capture to intercept computation
is approximately 15\,ms.

\subsection{Distillation}

The mathematical controller (Terms~I--III) is distilled into
a compact policy via MuJoCo rollouts:
\begin{enumerate}[nosep]
  \item generate $10^5$ randomised throws (varying $x_0$, $v_0$,
  object geometry);
  \item record successful catches as $(s_t, a_t)$ pairs;
  \item train a phase-conditioned MLP:
  $\pi_\theta\colon \R^{20} \to \R^7$
  (20-dim state $\to$ 7-dim joint torque);
  \item validate on held-out throw distribution.
\end{enumerate}

\begin{remark}[Distillation as renormalisation]
Compressing the continuous-time mathematical controller into a
discrete MLP is a renormalisation group step: the UV theory (full
state space, infinite precision, continuous time) flows to the IR
theory (finite network, finite precision, discrete control).
The spectral gap $\lambda_1 > 0$ is preserved through distillation
because the topological structure (cycle graph after capture) is
discrete and robust to approximation errors.  The $C_3$ contact
graph has $\lambda_1 = 3$, which is a macroscopic gap---no
continuous deformation of the policy can close it without breaking
a contact edge.
\end{remark}


% ═══════════════════════════════════════════════════════════
\section*{Closing: from Newton to the chopstick}
\addcontentsline{toc}{section}{Closing: from Newton to the chopstick}

Newton demonstrated that all objects fall at the same rate.
Galileo demonstrated it first, from the tower.  The equivalence
principle---inertial mass equals gravitational mass---is the
oldest universality in physics.  We have traced its consequence
through the full chain:

\begin{center}
\begin{tabular}{@{}rl@{}}
\textbf{Newton} & Projectile: $\ddot{x} = g$, mass-free. \\
\textbf{Equivalence} & The controller operates in $TM$, where
  $m$ has cancelled (\cref{thm:c-equivalence}). \\
\textbf{Sword} & The thrown object is an approaching sword
  (\cref{def:c-sword}). \\
\textbf{Well} & The chopstick pair creates a potential well
  (\cref{def:c-chopstick}). \\
\textbf{Capture} & Topological transition: $\lambda_1\colon 0
  \to 3$ (\cref{thm:c-capture}). \\
\textbf{Three terms} & Predict $+$ intercept $+$ bang
  (\cref{thm:c-decomp}). \\
\textbf{Kakeya} & No free capture: active control is necessary
  (\cref{cor:c-kakeya}). \\
\textbf{Implementation} & Flexiv $+$ chopsticks: the potential
  well made physical (\cref{sec:c-impl}).
\end{tabular}
\end{center}

Catching is the creation of a potential well at the right place
and time.  The sword enters the well, the well closes, and the
sword is no longer a sword.
